% ---------------------------------------------------------------- 页面
% 四种版式：screen / print / ereader / phone，见 tex/layout.tex。
\usepackage{geometry}
\ifbknarrow
  \ifx\booklayout\bkPhone
    \geometry{paperwidth=88mm, paperheight=176mm,
      left=6mm, right=6mm, top=6mm, bottom=8mm,
      headsep=0pt, footskip=5.5mm}
  \else
    \geometry{paperwidth=91mm, paperheight=122mm,
      left=5mm, right=5mm, top=5mm, bottom=7mm,
      headsep=0pt, footskip=5.5mm}
  \fi
\else
  \ifbkprint
    % 印刷版：内侧留装订线，页边不对称
    \geometry{a4paper, inner=33mm, outer=19mm, top=27mm, bottom=28mm,
      headsep=8mm, footskip=12mm}
  \else
    \geometry{a4paper, inner=27mm, outer=23mm, top=26mm, bottom=26mm,
      headsep=7mm, footskip=11mm}
  \fi
\fi

% ---------------------------------------------------------------- 字体
% 西文正文用 TeX Gyre Pagella（GUST 许可，可商用）。
\usepackage{fontspec}
\setmainfont{TeX Gyre Pagella}
% 中文与等宽字体。构建脚本会把 nixpkgs 里的字体路径写进 build/fonts.tex；
% 存在时按文件与 FontIndex 取精确字面，构建可复现，不依赖系统字体。
% 不存在时按字族名查找，最后退回 Fandol 与 lmmono，保证在没有 Nix 的机器上也能编译。
\IfFileExists{build/fonts.tex}{%
  \typeout{== book: 使用 build/fonts.tex 中的字体路径}%
  \input{build/fonts.tex}%
  \setCJKmainfont{SourceHanSerif.ttc}[
    Path=\shsdir,
    FontIndex=12,
    BoldFont=SourceHanSerif.ttc,
    BoldFeatures={FontIndex=27}
  ]
  \setCJKsansfont{Sarasa-Regular.ttc}[
    Path=\sarasadir,
    FontIndex=1,
    BoldFont=Sarasa-SemiBold.ttc,
    BoldFeatures={FontIndex=1}
  ]
  \setCJKmonofont{Sarasa-Regular.ttc}[
    Path=\sarasadir,
    FontIndex=37,
    BoldFont=Sarasa-SemiBold.ttc,
    BoldFeatures={FontIndex=37}
  ]
  \setsansfont{Sarasa-Regular.ttc}[
    Path=\sarasadir,
    FontIndex=1,
    BoldFont=Sarasa-SemiBold.ttc,
    BoldFeatures={FontIndex=1}
  ]
  \setmonofont{RecursiveMonoCslSt-Regular.otf}[
    Path=\recursivedir,
    Scale=0.95,
    BoldFont=RecursiveMonoCslSt-Bold.otf,
    ItalicFont=RecursiveMonoCslSt-Italic.otf,
    BoldItalicFont=RecursiveMonoCslSt-BdItalic.otf
  ]
}{%
  \PackageWarning{book}{未找到 build/fonts.tex，按字族名查找字体}%
  \IfFontExistsTF{Sarasa Fixed SC}{%
    \setmonofont{Sarasa Fixed SC}%
    \setCJKmonofont{Sarasa Fixed SC}%
    \setCJKsansfont{Sarasa Gothic SC}[BoldFont={Sarasa Gothic SC SemiBold}]%
    \setsansfont{Sarasa Gothic SC}[BoldFont={Sarasa Gothic SC SemiBold}]%
  }{%
    \setmonofont{Latin Modern Mono}%
    \setCJKmonofont{FandolHei-Regular.otf}%
    \setCJKsansfont{FandolHei-Regular.otf}[BoldFont=FandolHei-Bold.otf]%
    \setsansfont{Latin Modern Sans}%
  }%
  \IfFontExistsTF{Source Han Serif SC}{%
    \setCJKmainfont{Source Han Serif SC}[BoldFont={Source Han Serif SC Bold}]%
  }{%
    \setCJKmainfont{FandolSong-Regular.otf}[
      BoldFont=FandolSong-Bold.otf,
      ItalicFont=FandolKai-Regular.otf
    ]%
  }%
}
\newCJKfontfamily\cjkhei{Sarasa Gothic SC}[BoldFont={Sarasa Gothic SC SemiBold}]

% ---------------------------------------------------------------- 颜色
\usepackage{xcolor}
\definecolor{ink}{HTML}{1C3A5E}       % 主色：深墨蓝
\definecolor{muted}{HTML}{6B7280}     % 次要文字
\definecolor{rule}{HTML}{D9DEE4}      % 细线
\definecolor{paper}{HTML}{F5F6F7}     % 提示框底色
\definecolor{codebg}{HTML}{F7F8FA}    % 代码底色
% 插图用的第二层色：低饱和、可灰度打印、互相可区分
\definecolor{teal}{HTML}{2A6F6B}
\definecolor{ochre}{HTML}{A8742B}
\definecolor{clay}{HTML}{9A4B3F}
\definecolor{sage}{HTML}{5C7F58}
\definecolor{plum}{HTML}{6B4E71}
% 灰度版式：四种强调色换成四档明度，保证在电子墨水屏与黑白打印上可区分。
\ifbkgrayscale
  \definecolor{ink}{HTML}{1A1A1A}
  \definecolor{muted}{HTML}{5A5A5A}
  \definecolor{rule}{HTML}{C4C4C4}
  \definecolor{paper}{HTML}{F2F2F2}
  \definecolor{codebg}{HTML}{F5F5F5}
  \definecolor{teal}{HTML}{3C3C3C}
  \definecolor{ochre}{HTML}{6E6E6E}
  \definecolor{clay}{HTML}{4A4A4A}
  \definecolor{sage}{HTML}{8A8A8A}
  \definecolor{plum}{HTML}{2A2A2A}
\fi
% 手机是彩色屏，但版面只有 88mm 宽，图经 \figfit 缩放后图内标签会掉到
% 5pt 上下、结构线也只剩零点几毫米。次要文字与细线沿用灰度版式的深值，
% 否则这些元素在手机上近乎看不清。
\ifx\booklayout\bkPhone
  \definecolor{muted}{HTML}{5A5A5A}
  \definecolor{rule}{HTML}{8E949C}
\fi

% ---------------------------------------------------------------- 间距
\ifbknarrow
  % 正文 9pt、行距 11.5pt；再乘 1.28，实际基线距约 14.7pt。
  \linespread{1.28}
\else
  \linespread{1.38}
\fi
\ifbknarrow
  \setlength{\parindent}{1.6em}
\else
  \setlength{\parindent}{2em}
\fi
\setlength{\parskip}{0pt}
\setlength{\emergencystretch}{3em}
% 窄版式的行宽小，微小的溢出（<2pt）不可见；放宽 hfuzz 减少无意义的告警，
% 真正的溢出仍会报告。
\ifbknarrow
  \setlength{\hfuzz}{2pt}
  % 9pt 正文配 11.5pt 行距，接近 6 寸屏的舒适密度。
  \AtBeginDocument{\fontsize{9pt}{11.5pt}\selectfont}
  \renewcommand{\texttt}[1]{{\footnotesize\ttfamily #1}}
\fi
\raggedbottom
\clubpenalty=10000
\widowpenalty=10000
\displaywidowpenalty=10000

% ---------------------------------------------------------------- 标题
\usepackage{indentfirst}
\usepackage{titlesec}
\titleformat{\part}[display]
  {\raggedright}
  {\color{muted}\large\bfseries \CTEXthepart}
  {8pt}
  {\color{ink}\ifbknarrow\fontsize{17}{21}\else\fontsize{28}{34}\fi\selectfont\bfseries}
\ifbknarrow
  \titlespacing*{\part}{0pt}{18pt}{14pt}
\else
  \titlespacing*{\part}{0pt}{80pt}{40pt}
\fi

\titleformat{\chapter}[display]
  {\raggedright}
  {\color{muted}\normalsize\bfseries \CTEXthechapter}
  {6pt}
  {\color{ink}\ifbknarrow\fontsize{15}{19}\else\fontsize{24}{30}\fi\selectfont\bfseries}
  [\vspace{6pt}{\color{rule}\rule{\textwidth}{0.5pt}}]
\ifbknarrow
  \titlespacing*{\chapter}{0pt}{0pt}{12pt}
\else
  \titlespacing*{\chapter}{0pt}{0pt}{26pt}
\fi

\titleformat{\section}
  {\large\bfseries\color{ink}}
  {\thesection}{0.7em}{}
\titlespacing*{\section}{0pt}{20pt}{9pt}

% 小节标题不留在页底：后面至少留六行，否则整个标题移到下一页。
\makeatletter
\let\book@oldsection\section
\renewcommand{\section}{%
  \@ifstar{\book@oldsection*}{\bookneedspace{6\baselineskip}\book@oldsection}}
\makeatother

\titleformat{\subsection}
  {\normalsize\bfseries\color{ink!85}}
  {\thesubsection}{0.6em}{}
\titlespacing*{\subsection}{0pt}{15pt}{6pt}

% 章首导语
\newcommand{\chapterlead}[1]{%
  \begingroup
  \color{muted}\ifbknarrow\normalsize\else\large\fi\noindent
  #1\par
  \endgroup
  \ifbknarrow\vspace{8pt}\else\vspace{14pt}\fi%
}

% ---------------------------------------------------------------- 页眉页脚
\usepackage{fancyhdr}
\ifbknarrow\setlength{\headheight}{0pt}\else\setlength{\headheight}{14pt}\fi
\fancypagestyle{main}{%
  \fancyhf{}%
  \fancyhead[L]{\small\color{muted}\nouppercase{\leftmark}}%
  \fancyhead[R]{\small\color{muted}\nouppercase{\rightmark}}%
  \fancyfoot[C]{\small\color{muted}\thepage}%
  \renewcommand{\headrulewidth}{0pt}%
}
\fancypagestyle{plain}{%
  \fancyhf{}%
  \fancyfoot[C]{\small\color{muted}\thepage}%
  \renewcommand{\headrulewidth}{0pt}%
}
\renewcommand{\chaptermark}[1]{\markboth{\CTEXthechapter\quad #1}{}}
\renewcommand{\sectionmark}[1]{\markright{\thesection\quad #1}}
\ifbknarrow\pagestyle{plain}\else\pagestyle{main}\fi
\renewcommand{\headrulewidth}{0pt}

% 双面版式下补一张不编号的空白页，使下一节从右手页开始。
% \cleardoublepage 自动插入的空白页会带上页码，这里改用 \pagestyle{empty}。
\makeatletter
\newcommand{\bookblankverso}{%
  \clearpage
  \if@twoside\ifodd\c@page\else
    \pagestyle{empty}\null\newpage
    \ifbknarrow\pagestyle{plain}\else\pagestyle{main}\fi
  \fi\fi}
\makeatother

% ---------------------------------------------------------------- 目录
\usepackage{titletoc}
\titlecontents{part}[0pt]
  {\addvspace{16pt}\bfseries\color{ink}}
  {\thecontentslabel\quad}{}
  {\hfill\contentspage}
\titlecontents{chapter}[0pt]
  {\addvspace{8pt}\bfseries\color{ink}}
  {\thecontentslabel\quad}{}
  {\titlerule*[0.8pc]{.}\contentspage}
\titlecontents{section}[1.8em]
  {\addvspace{2pt}}
  {\thecontentslabel\quad}{}
  {\titlerule*[0.8pc]{.}\contentspage}

% ---------------------------------------------------------------- 表格
\usepackage{booktabs}
\usepackage{tabularx}
\usepackage{array}
\renewcommand{\arraystretch}{1.35}
% 列宽用比例而非绝对长度：表格总宽始终等于版心，换版式不会溢出。
% 参数是该列所占的份额，不必求和为整数；未标份额的 X 列份额为 1。
\newcolumntype{L}[1]{>{\raggedright\arraybackslash\hsize=#1\hsize}X}
\newcolumntype{C}[1]{>{\centering\arraybackslash\hsize=#1\hsize}X}
\newcolumntype{J}[1]{>{\hsize=#1\hsize}X}
% 表格允许跨页：窄页上任何一张表都可能高于版心，不可分割的 tabular 会造成无限分页。
\usepackage{etoolbox}
\usepackage{xltabular}
% 窄版式把表内字号降到 \scriptsize：88mm 与 91mm 宽的页上，9pt 的表格
% 一行常折成三行，整表比版心还高，跨页时机稍有不利就会触发无限分页。
% 钩子必须挂在 booktable 上：\AtBeginEnvironment 改写的是 \begin，
% 环境名不一致时不生效。
\newcommand{\bkttscriptsize}[1]{{\scriptsize\ttfamily #1}}
\AtBeginEnvironment{booktable}{%
  \ifbknarrow
    \scriptsize
    % 窄版式的 \texttt 用 \footnotesize（绝对字号），不会随 \scriptsize 变小，
    % 表内会出现代码比汉字还大的情况。表内改用 \let 换成一个 \scriptsize 版本：
    % 用 \renewcommand 带参数写法会在 etoolbox 的钩子里被再展开一次，实参丢失。
    \let\texttt\bkttscriptsize
  \fi}
% 表格一律用 xltabular（longtable 之上加 X 列）：可跨页，因此窄页上的长表
% 不会溢出页面，宽页上的短表也不会被强制分页。此前在宽页改用 tabularx 会
% 让按比例分配的 X 列多出约两倍行内边距，出现固定的 24pt 溢出。
\let\booktable\xltabular
\let\endbooktable\endxltabular

% ---------------------------------------------------------------- 列表
\usepackage{enumitem}
\ifbknarrow
  \setlist[itemize]{leftmargin=1.1em, itemsep=2pt, topsep=3pt, parsep=0pt}
  \setlist[enumerate]{leftmargin=1.4em, itemsep=2pt, topsep=3pt, parsep=0pt}
\else
  \setlist[itemize]{leftmargin=1.6em, itemsep=3pt, topsep=4pt, parsep=0pt}
  \setlist[enumerate]{leftmargin=1.8em, itemsep=3pt, topsep=4pt, parsep=0pt}
\fi

% ---------------------------------------------------------------- 图
\usepackage{float}
\ifbknarrow
  % 窄页上任何浮动体都可能超过版心高度，LaTeX 会因此不断生成空的浮动页。
  % 就地放置（H）不会推迟，也就不会循环。
  \floatplacement{figure}{H}
  \floatplacement{table}{H}
\else
  \floatplacement{figure}{htbp}
  \floatplacement{table}{htbp}
\fi
\usepackage{adjustbox}
\usepackage{tikz}
% 图形宽度保护：任何图都不超过版心，窄版式下自动等比缩小。
% 图同时限制宽与高：只限宽时，窄版式里高图会高于版心，[H] 浮动体放不下，
% 会在 \end{document} 处反复输出空页（表现为构建卡住、页数暴涨）。
\newcommand{\figfit}[1]{\adjustbox{max width=\linewidth, max height=0.72\textheight}{#1}}
\usetikzlibrary{arrows.meta, positioning, calc}
\usepackage{caption}
\captionsetup{font=small, labelfont={bf,color=ink}, labelsep=quad}

% ---------------------------------------------------------------- 代码
\usepackage{minted}
\usemintedstyle{tango}
% 代码块的底色盒需要内边距，否则首行贴住色块上沿、行首贴住左沿。
% minted 的 bgcolor 是 \fboxsep=0 的 \colorbox，本身没有留白；
% bgcolorpadding 会改用零宽边框（颜色同底色）撑出 framesep，肉眼只见留白。
\ifbknarrow
  \setminted{
    fontsize=\scriptsize,
    breaklines=true,
    breakanywhere=true,
    autogobble=true,
    bgcolor=codebg,
    bgcolorpadding=4pt,
    tabsize=2,
    xleftmargin=1.5pt,
    xrightmargin=1.5pt,
    baselinestretch=1.1
  }
\else
  \setminted{
    fontsize=\small,
    breaklines=true,
    breakanywhere=true,
    autogobble=true,
    bgcolor=codebg,
    bgcolorpadding=6pt,
    tabsize=2,
    xleftmargin=4pt,
    xrightmargin=4pt,
    baselinestretch=1.15
  }
\fi
\setmintedinline{bgcolor=codebg}

% 长标识符（包名、路径、选项名）在行内代码里无法自动换行，会在窄列或行尾
% 顶出版心。\bkbrk 在点、斜线、下划线与连字符之后插入可断点，只用于行内代码。
\newcommand{\bkbrk}[1]{#1\allowbreak}
\newcommand{\code}[1]{\texttt{#1}}
% shell 会话：$ 之后按 bash 高亮，其余行按普通文本输出，避免命令输出被染色。
\newminted[shellsession]{shell-session}{}
\newminted{nix}{}
\newminted{text}{}

% 代码块与表格不要贴着页底开始：剩余空间不足时整块移到下一页。
% 提示框是不可分割的整体，TeX 本来就会整块后移，无需预留。
% 不用 needspace 宏包：它靠一个 \penalty -100 制造断点，页面后段一旦溢出，
% TeX 会在代价比较中选中这个断点，即使当前页还剩七成空间，于是留下大片空白。
% 这里改成直接比较剩余高度，只在确实放不下时 \newpage。
\makeatletter
\newcommand{\bookneedspace}[1]{%
  \par
  \begingroup
    \dimen@=#1\relax
    \ifdim\pagegoal<\maxdimen
      \ifdim\dimen@>\pagegoal-\pagetotal\relax
        \newpage
      \fi
    \fi
  \endgroup}
\makeatother
\BeforeBeginEnvironment{shellsession}{\bookneedspace{5\baselineskip}}
\BeforeBeginEnvironment{nixcode}{\bookneedspace{9\baselineskip}}
\BeforeBeginEnvironment{textcode}{\bookneedspace{5\baselineskip}}
% 表格可跨页且会重复表头，只需保证表头与两行数据放得下，不要整表预留。
\BeforeBeginEnvironment{booktable}{\bookneedspace{5\baselineskip}}

% ---------------------------------------------------------------- 提示框
\usepackage[most]{tcolorbox}
% 提示框一律不可分割。允许跨页时，框首或框尾常单独留在页边，只剩标题加一行，
% 比整块移到下一页更难读。不可分割的代价是框前面的页尾留白，取舍见 docs/build.md。
\newcommand{\bkbreak}{}

\newtcolorbox{notebox}[1][提示]{%
  enhanced,
  colback=paper, colframe=ink!55,
  boxrule=0pt, leftrule=2pt, arc=0pt, outer arc=0pt,
  left=10pt, right=10pt, top=8pt, bottom=8pt,
  before skip=12pt, after skip=12pt,
  before upper={\bfseries\small\color{ink}#1\par\vspace{3pt}}%
}
\newtcolorbox{warnbox}[1][注意]{%
  enhanced,
  colback=paper, colframe=ink!85,
  boxrule=0pt, leftrule=2pt, arc=0pt, outer arc=0pt,
  left=10pt, right=10pt, top=8pt, bottom=8pt,
  before skip=12pt, after skip=12pt,
  before upper={\bfseries\small\color{ink}#1\par\vspace{3pt}}%
}
% 检查点框在窄版式下允许跨页：条目多的章（八条以上）框高会超过版心，
% 不可分割的框一旦高于版心，TeX 会在页首反复重排，生成上万页空页。
% 宽版式仍不可分割，保持整块外观。
\ifbknarrow
  \def\bkcheckbreak{breakable}
\else
  \def\bkcheckbreak{}
\fi
\newtcolorbox{checkpoint}[1][本章检查点]{%
  enhanced, \bkcheckbreak,
  colback=white, colframe=rule,
  boxrule=0.5pt, arc=0pt, outer arc=0pt,
  left=12pt, right=12pt, top=10pt, bottom=10pt,
  before skip=16pt, after skip=12pt,
  before upper={\bfseries\small\color{ink}#1\par\vspace{4pt}}%
}
\newtcolorbox{lab}[1][十分钟上手]{%
  enhanced,
  colback=white, colframe=teal!70,
  boxrule=0pt, leftrule=2.5pt, arc=0pt, outer arc=0pt,
  left=10pt, right=10pt, top=7pt, bottom=7pt,
  before skip=12pt, after skip=12pt,
  before upper={\bfseries\small\color{teal}#1\par\vspace{4pt}}%
}
\newtcolorbox{culturenote}{%
  enhanced,
  colback=ochre!5, colframe=ochre!70,
  boxrule=0pt, leftrule=2pt, arc=0pt, outer arc=0pt,
  left=10pt, right=10pt, top=7pt, bottom=7pt,
  before skip=12pt, after skip=12pt,
  before upper={\bfseries\small\color{ochre}Unix 文化\par\vspace{3pt}}%
}
\newtcolorbox{toolswap}[2]{%
  enhanced,
  colback=paper, colframe=rule,
  boxrule=0pt, leftrule=2pt, arc=0pt, outer arc=0pt,
  left=10pt, right=10pt, top=7pt, bottom=7pt,
  before skip=12pt, after skip=12pt,
  before upper={\bfseries\small\color{ink}旧 \textrightarrow\ 新\quad
    \texttt{#1}\ \textrightarrow\ \texttt{#2}\par\vspace{3pt}}%
}

% ---------------------------------------------------------------- 链接
\usepackage{hyperref}
\hypersetup{
  hidelinks,
  pdftitle={\booktitle：\booksubtitle},
  pdfauthor={\bookauthor},
  pdfsubject={Linux 与 NixOS 入门},
  pdfkeywords={Linux, NixOS, 命令行, 服务器运维, AI Agent, \bookversion}
}
\usepackage{url}
\urlstyle{same}
